\begin{abstract}
针对六区域异构 AI 任务与新能源、储能、电网的耦合调度，本文构建“闭环需求预测—分层任务邻域—离散 SOC 动态规划—按情景交替搜索”的可复算框架。问题一在独立验证集上从季节朴素与 Huber 候选中选模，18 条序列中 17 条选择 Huber，平均 MAE、RMSE 和 WAPE 分别为 25.5892 GPU、34.5070 GPU 和 0.8231。问题二每个方向评价 10000 个候选；五维分层使单方案覆盖率达到 3.038\%--3.616\%，合计覆盖 16.148\%；平衡方案成本为 $-4.213783\times10^8$ 元，碳排放 8.7569 tCO$_2$。问题三对不重复非支配方案统一采用加权理想点距离，选择成本策略，成本由 $-4.198203\times10^8$ 元降至 $-4.587522\times10^8$ 元，碳排放降至 0。问题四的联合基准与 17 个情景统一使用 21 个 SOC 状态；联合成本、碳排放和峰值购电分别为 $-4.587886\times10^8$ 元、0 tCO$_2$ 和 0 MW。两次独立运行的结果哈希和 Q1--Q4 语义一致。全部结论均为有限搜索 best-found 的场景可行结果，不构成全局最优或完整 Pareto 前沿。
\end{abstract}
\keywords{异构 AI 任务；闭环预测；算电协同；储能动态规划；多目标调度}
